\appendix

%\renewcommand{\thesection}{A}
%\setcounter{equation}{0}
%\setcounter{theor}{0}
%\addtocontents{toc}{\protect\vspace{-2ex}}

%\addcontentsline{toc}{section}{Appendix A. The hyperbolic gamma function}
																																												 
%\section*{Appendix A. The hyperbolic gamma function}
\section{The hyperbolic gamma function}\label{appendixA}

The  hyperbolic gamma function was introduced and studied in~\cite{fo} as a so-called minimal solution of a special f\/irst order analytic dif\/ference equation. It is basically the same as Kurokawa's double sine~\cite{kuro}, Faddeev's quantum dilogarithm~\cite{fadd2}, and Woronowicz's quantum exponential function~\cite{woro}. (The precise connections between these functions are spelled out in Appendix~A of our paper~\cite{Hels}.) In this appendix we review features of the hyperbolic gamma function $G(\aaa;z)$ that are used in the present paper; if need be, see~\cite{fo} for proofs.

Unless specif\/ied otherwise, we choose
\begin{gather}
\aaa >0,
\end{gather}
and suppress the dependence of $G$ on $a_+$, $a_-$.

 To begin with, $G(z)$ can be def\/ined as the unique minimal solution of one of the two analytic dif\/ference equations
\begin{equation}\label{Gades}
	\frac{G(z+ia_\delta/2)}{G(z-ia_\delta/2)} = 2c_{-\delta}(z),\qquad  \delta=+,-,
\end{equation}
that has modulus 1 for real $z$ and satisf\/ies $G(0)=1$ (recall~\eqref{denot} for the notation used here); remarkably, this entails that the other one is then satisf\/ied as well.
It is meromorphic in~$z$,
and for $z$ in the strip
\begin{gather}\label{strip}
S\equiv \{z\in\C \mid |\im (z)|<a\},
\end{gather}
no poles and zeros occur. Hence we have
\begin{gather}\label{Gg}
G(z)=\exp(ig(z)),\qquad z\in S,
\end{gather}
with the function $g(z)$ being holomorphic in $S$.
 Explicitly, $g(z)$ has the integral representation
\begin{gather}\label{ghyp}
g(a_{+},a_{-};z) =\int_0^\infty\frac{dy}{y}\left(\frac{\sin 2yz}{2\sinh(a_{+}y)\sinh(a_{-}y)} - \frac{z}{a_{+}a_{-} y}\right),\qquad z\in S.
\end{gather}
From this, the following properties of the hyperbolic gamma function are immediate:
\begin{gather}\label{refl}	
G(-z) = 1/G(z),\qquad ({\rm ref\/lection\ equation}),
\\
\label{modinv}
G(a_-,a_+;z) = G(a_+,a_-;z),\qquad  ({\rm modular\ invariance}),
\\
\label{sc}
	G(\lambda a_+,\lambda a_-;\lambda z) = G(a_+,a_-;z),\qquad \lambda\in(0,\infty),\qquad ({\rm scale\ invariance}),
\\
\label{Gcon}
\overline{G(\aaa;z)}=G(\aaa;-\overline{z}).
\end{gather}

We have occasion to use a few more features that are less obvious, including the duplication formula
\begin{gather}\label{dupl}
G(\aaa;2z)=G(\aaa;z\pm ia_{+}/4 \pm ia_{-}/4),
\end{gather}
the closely related formula
\begin{gather}\label{G2}
G(a_{+},2a_{-};2z)=G(\aaa;z\pm ia_{+}/4),
\end{gather}
the explicit evaluation
\begin{gather}\label{Geval}
G(ia_{+}/2-ia_{-}/2)=(a_{+}/a_{-})^{1/2},
\end{gather}
and the asymptotic behavior of $G(z)$ for $\re (z)\to\pm \infty$. The latter is given by
\begin{gather}\label{Gas}
	G(\aaa;z) = \exp \big({\mp} i\left(\chi+\alpha z^2/4\right)\big)\big(1 + O(\exp(-r |\re(z)|))\big),\qquad \re(z)\to\pm\infty,
\end{gather}
where the decay rate $r$ can be any positive number satisfying
\begin{gather}
r <\alpha \min(a_+,a_-),
\end{gather}
and where
\begin{equation}\label{chi}
 \chi \equiv \frac{\pi}{24}\left(\frac{a_+}{a_-} + \frac{a_-}{a_+}\right).
\end{equation}

Def\/ining
\begin{gather}\label{zkl}
z_{kl}\equiv ika_{+}+ila_{-},\qquad k,l\in\N\equiv \{ 0,1,2,\ldots\},
\end{gather}
the hyperbolic gamma function has its poles at
\begin{gather}\label{Gpo}
z=z_{kl}^-,\qquad z_{kl}^-\equiv -ia -z_{kl},\qquad k,l\in\N,\qquad (G\mbox{-poles}),
\end{gather}
and its zeros at
\begin{gather}\label{Gze}
z=z_{kl}^+,\qquad z_{kl}^+\equiv ia +z_{kl},\qquad k,l\in\N,\qquad (G\mbox{-zeros}).
\end{gather}
The pole at $-ia$ is simple and has residue
\begin{gather}\label{Gres}
\lim_{z\to -ia}(z+ia)G(z)=\frac{i}{2\pi}(a_{+}a_{-})^{1/2}.
\end{gather}
In view of these features, $G(z)$ can be written as a ratio of entire functions,
\begin{gather}\label{GE}
G(\aaa;z)=E(\aaa;z)/E(\aaa;-z),
\end{gather}
where $E(\aaa;z)$ has its zeros at
\begin{gather}\label{Eze}
z=z^{+}_{kl},\qquad k,l\in\N,\qquad (E\mbox{-zeros}).
\end{gather}

The function $E(\aaa;z)$ we have occasion to employ is def\/ined in Appendix~A of~I; it is closely related to Barnes' double gamma function~\cite{barn}. We need two more of its properties. First, from equations~(A.41) and~(A.43) in~I we have
 \begin{gather}\label{EwT}
 E(z)E(-z)=\exp\left(\frac12\int_0^{\infty}\!\frac{dy}{y}\left(
 \frac{1-\cos(2yz)}{\sinh(a_{+}y)\sinh(a_{-}y)}-\frac{z^2}{a_+a_-}\big(e^{-2a_{+}y}+e^{-2a_{-}y}\big)\right)\right),\!\!\!
 \end{gather}
where $z$ belongs to the strip $S$~\eqref{strip}. Second, we need the A$\De$Es it satisf\/ies, namely,
\begin{gather}\label{Eades}
\frac{E(z+ia_{\de}/2)}{E(z-ia_{\de}/2)}=\sqrt{2\pi}\exp(izK_{\de})/\Gamma(iz/a_{-\de}+1/2),\qquad \de=+,-,
\\
K_{\de}\equiv \frac{1}{2a_{-\de}}\ln \left( \frac{a_{\de}}{a_{-\de}}\right),
\end{gather}
cf.~equations~(A.46)--(A.47) in~I.

We also state two zero step size limits of the hyperbolic gamma function, which we need for taking nonrelativistic limits. The f\/irst one yields the relation to the Euler gamma function:
\begin{gather}\label{GhGr}
\lim_{\kappa\downarrow 0}G(1,\kappa;
 \kappa z+i/2)\exp \big(iz\ln(2\pi\kappa)-\ln(2\pi)/2\big) = 1/\Gamma(iz+1/2).
\end{gather}
For the second one we need to require that $z$ stay away from  cuts given by $\pm i[a_+/2,\infty)$. Then we have
\begin{gather}\label{Gnr}
\lim_{a_{-}\downarrow 0}\frac{G(a_+, a_{-};z+iua_{-})}{G(a_+, a_{-};z+ida_{-})}
=\exp((u-d)\ln [2c_+(z)]),\qquad u,d\in\R,
\end{gather}
